% Generated by manuscript2/tables2/make_prnp_somatic_snv_summary_tex.py
\begin{table}[t]
  \centering
  \resizebox{\textwidth}{!}{
    \begin{tabular}{lllcccccrrrrrrrrr}
      \toprule
      Sample & Chromosome & Position & VAF (\%) & REF & ALT & Mutation type
      & Population frequency & Read depth & REF count & ALT count & Mean BQ & Mean MQ
      & REF fwd & REF rev & ALT fwd & ALT rev \\
      \midrule
      CJD23 & Chr20 & 4691920 & 0.95 & G & A & intron & not described & 5178 & 5129 & 49 & 21.00 & 38.49 & 2538 & 2591 & 17 & 32 \\
      CJD2 & Chr20 & 4691920 & 0.89 & G & A & intron & not described & 4738 & 4696 & 42 & 22.01 & 37.90 & 2336 & 2360 & 14 & 28 \\
      CJD23 & Chr20 & 4694249 & 0.82 & T & C & intron & not described & 1346 & 1335 & 11 & 20.62 & 30.01 & 687 & 648 & 5 & 6 \\
      \bottomrule
    \end{tabular}
  }
  \vspace{0.5em}
  \caption{\textbf{PRNP intronic calls retained by sequencing2.} Three sample-level calls,
  representing two genomic sites, passed all configured filters, including the empirical
  complete-pipeline LoD of 26/3891 (0.668\%). VAFs ranged from 0.82\% to 0.95\%.
  Abbreviations: VAF, variant allele fraction; REF, reference allele; ALT, alternative allele;
  BQ, base quality; MQ, mapping quality; fwd, forward; rev, reverse.}
  \label{tab:sequencing2_prnp_somatic_snv_summary}
\end{table}

\clearpage
